%MIT OpenCourseWare: https://ocw.mit.edu
%RES.18-012 Algebra II Student Notes, Spring 2022
%License: Creative Commons BY-NC-SA 
%For information about citing these materials or our Terms of Use, visit: https://ocw.mit.edu/terms.

% \rhead{\rightmark}
\lhead{Lecture 14: Number Fields}
\section{Number Fields}

\subsection{The Gaussian Integers}

Last time, we discussed factorization in the Gaussian integers $\ZZ[i]$, and how it relates to the sum of squares problem. At the end of the lecture, we classified all Gaussian primes. Note that since $\ZZ[i]$ is a UFD, the primes in $\ZZ[i]$ are precisely the irreducible elements. In general, primes are defined as elements such that if $p \mid ab$, then $p \mid a$ or $p \mid b$ --- but in a UFD, an element is prime if and only if it's irreducible.

% Last time, we talked about factorization in $\ZZ[i],$ the Gaussian integers, and sums of squares. At the end of the lecture, we classified all the Gauss primes. In particular, since $\ZZ[i]$ is a UFD, the primes in $\ZZ[i]$ are precisely the irreducible elements.

% \begin{qq}
% What are the Gauss primes in $\ZZ[i]$? Which integers $n \in \ZZ$ can be written as a sum of squares?
% \end{qq}

\begin{theorem}
The complete list of all primes in $\ZZ[i]$, up to association, consists of:
\begin{itemize}
    \item for each integer prime $p = 4k + 3 \in \ZZ$, the Gaussian prime $p$ itself;
    \item for each integer prime $p = 4k + 1 \in \ZZ$, the two Gaussian primes $a \pm bi$ where $a^2 + b^2 = p$;
    \item the prime $1 + i$. 
\end{itemize}
\end{theorem}

% \begin{proof}
Note that we have $2 = (1 + i)(1 - i)$, but $1 + i$ and $1 - i$ are associate, so $2$ only contributes \emph{one} prime up to association. 

\begin{proof}
    First we'll check that all such elements are primes. For the second and third cases, where $p$ factors as $a^2 + b^2$ for some integers $a$ and $b$, it's enough to prove the following claim:

    \begin{claim*}
        If $a^2 + b^2 = p$ is an integer prime, then $a + bi$ is prime in $\ZZ[i]$. 
    \end{claim*}

    \begin{proof}
        Define the norm $\operatorname{N}(a + bi) = a^2 + b^2$, which is multiplicative. Suppose $a + bi$ factors as $\alpha\beta$. Then we have \[p = \operatorname{N}(a + bi) = \operatorname{N}(\alpha)\operatorname{N}(\beta).\] Then since $p$ is an integer prime, $\operatorname{N}(\alpha)$ or $\operatorname{N}(\beta)$ must be $1$. But if $\operatorname{N}(\alpha) = 1$, then $\alpha\overline{\alpha} = 1$, so $\alpha$ is a unit. This means in any factorization of $a + bi$, one factor must be a unit, so $a + bi$ is irreducible and therefore prime. 
    \end{proof}

    Meanwhile, for the first case, we saw last class that every integer prime $p = 4k + 3$ is still prime in $\ZZ[i]$. 

    Now we will check that there are no other primes --- it's enough to check that every non-unit $\alpha \in \ZZ[i]$ is divisible by some element of this list. To do so, we again use the norm --- if $\alpha$ is not a unit, then we have $\alpha\overline{\alpha} = n$ for some integer $n > 1$. Let $p$ be a prime divisor of $n$. 
    
    Then if $p \equiv 3 \pmod{4}$, we must have $p \mid \alpha$ (or $p \mid \overline{\alpha}$, which also implies $p \mid \alpha$), since $p$ is prime. Otherwise, we can write $p = (a + bi)(a - bi)$ where $a \pm bi$ are both primes. Then $a + bi$ must divide $\alpha$ or $\overline{\alpha}$, and therefore $a + bi$ or $a - bi$ must divide $\alpha$. 
\end{proof}

\begin{question}
Does this still work when $p = 2$?
\end{question}

\begin{ans}
Yes, the argument still works as written --- in fact, we don't even need the last step, since if $1 + i$ divides $\overline{\alpha}$, then $1 + i$ \emph{itself} also divides $\alpha$. 
\end{ans}

% If $a^2 + b^2 = p,$ a usual prime, then $a + bi \in \ZZ[i]$ is a prime in $\ZZ[i]$, which is equivalent to being irreducible. The norm is $N(a + bi) = a^2 + b^2.$ If there is a factorization $a + bi = \alpha \beta$, then $N(a + bi) = N(\alpha)N(\beta).$ Since $N(a + bi)$ is a usual prime, either $N(\alpha)$ or $N(\beta)$ is 1, so either $\alpha\overline{\alpha} = 1$ or $\beta\overline{\beta} = 1$, and then $\alpha$ or $\beta$ must be a unit. Therefore, the only divisors are units and associates, so $a + bi$ is irreducible and thus a prime. It is also the case that $p = 4k + 3$ is a prime in $\ZZ[i],$ and so all the listed elements are prime.

% Moreover, no other primes exist. To show this, it is enough to check that if $\alpha \in \ZZ[i]$ is not a unit, then it is divisible by an element on the list. If $\alpha$ is not a unit, then $\alpha \overline{\alpha} = n \neq 1.$ Let $p$ be a prime divisor of $n.$ If $p = 4k + 3,$ then $p$ is a prime, so $p$ divides either $\alpha$ or $\overline{\alpha},$ but if $p$ divides $\overline{\alpha},$ by conjugating, $p$ also divides $\alpha.$ If $p = 4k + 1$ or $p = 2,$ then $p = (a + bi)(a - bi),$ where $(a + bi)$ divides $\alpha$ or $\overline{\alpha}$. Similarly, if $a + bi$ divides $\overline{\alpha}$, then $a - bi$ divides $\alpha.$ In every case, $\alpha$ is divisible by some element on the list. Therefore, the list consists precisely of all the primes in $\ZZ[i].$ 
% % \end{proof}

As a corollary, we can find the complete answer to the sum of squares question. 

\begin{corollary}
If $n$ has prime factorization $n = p_1^{d_1}\cdots p_r^{d_r}$ in $\ZZ$, then $n$ is a sum of squares if and only if the exponent $d_i$ is even for all primes $p_i \equiv 3 \pmod{4}$. 
\end{corollary}

For example, $21$ has odd exponents of $3$ and $7$, so it cannot be written as a sum of squares. 

\begin{proof}
First, to show that all $n$ of this form work, we've seen earlier that if $m$ and $n$ are sums of squares, then so is $mn$. For all $p \not\equiv 3 \pmod{4}$, we've seen that $p$ is a sum of squares; and for all $p \equiv 3 \pmod{4}$, we have that $p^2$ is trivially a sum of squares. Since $n$ is the product of such terms, it must be a sum of squares as well. 
% The product of sums of squares is a sum of squares. Consider $n$ of this form. Then it is a product of primes $p$ that are not 3 modulo 4, which are sums of squares, and squares $m^2$ of primes that are 3 modulo 4, which are trivially sums of squares. So $n$ of this form is a sum of squares. 

Conversely, suppose $n$ can be written as a sum of squares, so $n = (a + bi)(a - bi)$. Let $d$ be the power of a given prime $p \equiv 3 \pmod{4}$ in the prime factorization of $a + bi$ in the Gaussian integers. Since $p \in \ZZ$, then $d$ is also the power of $p$ in the prime factorization of $a - bi$, so the power of $p$ in the factorization of $n$ is $2d$. Therefore, the power of each prime $p \equiv 3 \pmod{4}$ in the prime factorization of $n$ must be even.  
\end{proof}

\begin{question}
Why doesn't this statement have a condition involving $p = 2$?
\end{question}

\begin{ans}
This is because $2$ is a sum of squares, as $2 = 1 + 1$. So $2$ is allowed to have either odd or even power in the prime factorization of $n$. 
\end{ans}

This is just one example of how such considerations can be applied to number theory. It is possible to go even further --- for example, it is possible to determine how many different presentations as a sum of squares there are for a given $n$. 

\subsection{Fermat's Last Theorem, as an Aside}

The ideas used to analyze solutions to $n = a^2 + b^2$ have some relevance to a more difficult equation as well, Fermat's Last Theorem. 

% Consider the more difficult equation $a^n + b^n = c^n$ (over the integers).

\begin{theorem}[Fermat]
For an integer $n > 2$, the equation \[a^n + b^n = c^n\] has no solutions where $a$, $b$, and $c$ are all nonzero integers.  
\end{theorem}

Mathematicians have been trying to prove this theorem for a long time. Fermat famously proposed the theorem in the margin of a book, stating, ``I have a truly marvelous demonstration of this proposition that this margin is too narrow to contain.'' He most likely did not have a truly marvelous demonstration of the proposition. 

Mathematician Gabriel Lam\'e announced a proof on March $1$ of $1847$. This proof was later found to be incorrect, then partially corrected to be valid in certain cases by Ernst Kummer.

The initial steps of Lam\'e's proof proceed in a similar fashion as our analysis for Gaussian integers. Assume $n$ is odd (in fact, we can assume that $n$ is prime --- it suffices to prove the theorem for $n = 4$ and $n$ prime, and Fermat did have a proof for the case $n = 4$). When considering sums of squares, we used the factorization $a^2 + b^2 = (a + bi)(a - bi)$. In general, when $n$ is odd we have a similar factorization \[a^n + b^n = (a + b)(a + \zeta b)(a + \zeta^2 b) \cdots (a + \zeta^{n - 1}b),\] where $\zeta = e^{2\pi i/n}$. This gives a factorization of $a^n + b^n$ in the ring $\ZZ[\zeta]$, the ring of \emph{cyclotomic integers}.  

In many cases, it is possible to check that the factors are pairwise coprime; that is, that they do not have a common divisor. In the usual integers, if an $n$th power is factored as a product of coprime factors, then every factor must be also an $n$th power (up to multiplication by $\pm 1$). In this case, we similarly want to conclude that each factor is an $n$th power (up to multiplication by units). This would eventually lead to a contradiction. %This leads to a contradiction. 

The key point is that if $\ZZ[\zeta]$ were a UFD, then we could obtain our conclusion. This is what Lam\'e didn't show --- we're so used to dealing with the integers, where unique factorization \emph{does} hold, that even a major mathematician initially missed that unique factorization doesn't have to hold in this setting (though it was realized later). Unfortunately, this method doesn't work --- when $n$ is an odd prime, $\ZZ[\zeta]$ is \emph{almost never} a UFD. 

However, Kummer showed that under a weaker condition than $\ZZ[\zeta]$ being a UFD, it's still possible to obtain this conclusion. This weaker condition is that $p$ is a \emph{regular} prime --- not every prime is regular, but many are. In order to explain what a regular prime is, we need more theory, which we'll see in future classes. As a glimpse into what this theory will be, when unique factorization fails, we can still analyze \emph{how much} it fails. This leads to the definition of the \emph{ideal class group}, which in some sense controls the non-uniqueness of prime factorization; then the regularity of a prime is a property of the ideal class group of $\ZZ[\zeta]$. 

% The condition is that $p$ is a \emph{regular} prime --- this is a property of something that controls the non-uniqueness of factorization. (If unique factorization fails, we can still analyze \emph{how much} it fails.) More precisely, it's a property of the \emph{ideal class group}, which we'll see in a later class. 

% , which Lam\'e did not show at the time, is that if $\ZZ[\zeta]$ were a UFD, we could obtain our conclusion. Unfortunately, without loss of generality we can assume $n$ is a prime, in which case $\ZZ[\zeta]$ is never a UFD. However, under some weaker condition, the conclusion can still be reached. 

\subsection{Number Fields}

Now we will move on to a more general case. 

\begin{qq}
How does factorization work in rings like $\ZZ[i]$ and $\ZZ[\zeta]$?
\end{qq}

Both rings sit inside a larger field: $\ZZ[i] \subset \QQ[i]$ and $\ZZ[\zeta] \subset \QQ[\zeta]$. This is helpful because fields can be easier to work with than rings. The concept of a \emph{number field} generalizes this.

\begin{definition}
A \textbf{number field} is a subfield in $\CC$ that is finite-dimensional as a vector space over $\QQ$.
\end{definition}

\begin{example}
The number field $\QQ[i] = \{a + bi \mid a, b \in \QQ\}$ is two-dimensional over $\QQ$.
\end{example}

\subsubsection{Algebraic Numbers and Integers}

An important observation is that every element in a number field is algebraic.

\begin{definition}
A number is an \textbf{algebraic number} if it is a root of a polynomial with rational coefficients.
\end{definition}

If $F$ is a number field, and $\alpha \in F$, then there is a linear dependence between $1$, $\alpha$, $\alpha^2$, \ldots, $\alpha^n$ for any $n \geq \dim_{\QQ}(F)$. So $\alpha$ is a root of some polynomial $P \in \QQ[x]$, and therefore $\alpha$ is algebraic. Conversely, $\QQ[\alpha]$ is a number field if $\alpha$ is algebraic --- if $\alpha$ is the root of a polynomial of degree $d$, then we can express powers $\alpha^i$ with $i \geq d$ in terms of lower powers, so the only possible terms we have are $1$, $\alpha$, \ldots, $\alpha^{d - 1}$. This argument generalizes to show that $\QQ[\alpha_1, \ldots, \alpha_n]$ is also a number field if $\alpha_1$, \ldots, $\alpha_n$ are all algebraic, since there are only finitely many terms $\alpha_1^{e_1}\cdots \alpha_n^{e_n}$ which we need to consider (if a term contains a high power of some $\alpha_i$, we can rewrite it in terms of lower powers). 

If $\alpha$ is algebraic, then $\{P \mid P(\alpha) = 0\}$ is an ideal in $\QQ[x]$. Since $\QQ[x]$ is a PID, it has the form $(P)$, where $P$ is a monic polynomial of minimal degree. Such a $P$ is called the \textbf{minimal polynomial} for $\alpha$ over $\QQ$.

A number field is a generalization of fields like $\QQ[i]$ or $\QQ[\zeta]$, but we really want to analyze factorization in \emph{rings} like $\ZZ[i]$ or $\ZZ[\zeta]$, not the underlying fields. To describe the generalization of such rings to an arbitrary number field, we need to define an \emph{algebraic integer}. 

\begin{definition}
An algebraic number is an \textbf{algebraic integer} if its minimal polynomial has integer coefficients.
\end{definition}

\begin{lemma}
The element $\alpha$ is an algebraic integer if and only if $P(\alpha) = 0$ for \emph{some} monic polynomial $P \in \ZZ[x].$ 
\end{lemma}

It is evident that a polynomial with rational coefficients can be scaled to either be monic or have integer coefficients, but an algebraic integer requires that both can be achieved simultaneously.

\begin{proof}
The proof is another application of Gauss's Lemma and the ideas from last lecture. One direction is obvious: if the minimal polynomial $P$ is in $\ZZ[x]$, then $P(\alpha) = 0$, so the condition is clearly satisfied. 

For the other direction, suppose there exists a monic polynomial $P \in \ZZ[x]$ such that $P(\alpha) = 0$. Now consider the minimal polynomial $P_{\min}$. By clearing denominators and pulling out the gcd of its coefficients, rescale it to a polynomial $Q = aP_{\min}/b$ which is primitive and in $\ZZ[x]$. 

Then $Q$ divides $P$ in $\QQ[x]$, since $P_{\min}$ divides $P$. But $Q$ is primitive, so by the results in the last lecture, $Q$ must also divide $P$ in $\ZZ[x]$. But then the leading coefficient of $Q$ must divide the leading coefficient of $P$. Since $P$ is monic, $\pm Q$ must be monic as well. Then since $P_{\min}$ is also monic, we have $Q = \pm P_{\min}$, so $P_{\min}$ has integer coefficients. 
\end{proof}


\begin{example}
A rational number $\alpha \in \QQ$ has minimal polynomial $x - \alpha$, so $\alpha$ is be an algebraic integer if and only if $\alpha$ is a usual integer. 
\end{example}

The next example is the primary setting that we will work with.

\begin{example}[Quadratic Number Fields]
What are the algebraic integers in the number field $\QQ[\sqrt{d}]$? (Here $d$ may or may not be positive.)
\end{example}

\begin{proof}
Without loss of generality, we can assume $d$ is squarefree (since factoring squares out of $d$ doesn't change the number field). Then let $\alpha = a + b\sqrt{d}$ with $a, b \in \QQ$. The minimal polynomial of $\alpha$ is \[(x - a - b\sqrt{d})(x - a + b\sqrt{d}) = x^2 - 2a + (a - b^2d),\] so $\alpha$ is an algebraic integer if and only if $2a \in \ZZ$ and $a^2 - b^2d \in \ZZ$. 

Now we have two cases: 

\emph{Case 1 ($a \in \ZZ$).} Then we must have $b^2d \in \ZZ$ as well, and since $d$ is squarefree, $b \in \ZZ$ as well. 

\emph{Case 2 ($a = k + \frac{1}{2}$ for $k \in \ZZ$).} Then we must have $2b \in \ZZ$ as well, and $b = m + \frac{1}{2}$. In this case, we have \[a^2 - b^2d = \frac{1}{4}((2k + 1)^2 - (2m + 1)^2d).\] This is an integer if and only if $d \equiv 1 \pmod{4}$. 

So the conclusion is that if $d \not\equiv 1 \pmod{4}$, the algebraic integers are precisely $a + b\sqrt{d}$ for $a, b \in \ZZ$ --- for example, the algebraic integers in $\QQ[i]$ and $\QQ[\sqrt{-5}]$ are $\ZZ[i]$ and $\ZZ[\sqrt{-5}]$. Meanwhile, if $d \equiv 1 \pmod{4}$, the algebraic integers are precisely $a + b\sqrt{d}$ where $a, b \in \ZZ$ or $a + \frac{1}{2}, b + \frac{1}{2} \in \ZZ$ --- for example, the algebraic integers in $\QQ[\sqrt{-3}]$ are $\ZZ[\omega]$ for a primitive third root of unity $\omega$, which are exactly the elements of this form. 
\end{proof}

% \begin{example}[Quadratic Number Fields]
% Consider $F = \QQ[\sqrt{d}],$ where $d$ may or may not be positive. Without loss of generality, $d$ is squarefree.\footnote{A square part of $d$ would not change the field $\QQ$} When $\alpha = a + b \sqrt{d}$, the minimal polynomial is $x^2 - 2a + (a^2 - b^2d),$ so $\alpha$ is an algebraic integer if and only if $2a \in \ZZ$ and $a^2 - b^2 d$ is an integer. 
% \end{example}
% This example splits into two cases. If $a \in \ZZ,$ then $b^2 d \in \ZZ$ as well, and since $d$ is squarefree, $b \in \ZZ$ as well. The other case is that $a$ is a half-integer $k + 1/2$ for $k \in \ZZ$. Then $2b \in \ZZ$ as well, and $b = m + \frac{1}{2}.$ In this case, the polynomial becomes $\frac{1}{4}((2k + 1)^2 - (2m + 1)^2d) \in \ZZ$ and $d = 1$ modulo 4.

% The conclusion is that $\alpha = a + b \sqrt{d}$ where for $d \neq 1$ mod 4, $a$ and $b$ are integers, such as in $\ZZ[i]$ and $\ZZ[\sqrt{-5}],$ and for $d = 1$ mod 4, $a$ and $b$ are half-integers, which is true in fields like $\QQ[\sqrt{-3}].$

In general, the algebraic integers of a number field have quite a bit of additional structure. 

\begin{theorem}
For a number field $F$, the set of algebraic integers in $F$ is a subring of $F$. Furthermore, this is the largest subring that is finitely generated as an abelian group under addition. 
% The set of algebraic integers in a number field has quite a bit of additional structure.
% \begin{enumerate}[label = (\alph*)]
%     \item If $F$ is a number field, then the set of algebraic integers is a subring of $F$.
%     \item In fact, this is the largest subring that is finitely generated as an abelian group under addition.
% \end{enumerate}
\end{theorem}

\todo{post proof of this as part of the appendix??}

We won't prove this in general; but it's fairly straightforward in our specific example $F = \QQ[\sqrt{d}]$.

Next time, we will study factorization in rings of algebraic integers. The general idea is that if a statement fails to be true, mathematicians often try to generalize it to find a version that \emph{is} true. In this case, where unique factorization fails for a general ring of algebraic integers, we can consider a more general version of it which works. 

To motivate what this more general version will be, note that in the cases where unique factorization works, primes are considered up to association. In fact, an element in a ring $\alpha \in R$ up to association is equivalent to the ideal $(\alpha)$. So in the more general case, we actually consider products of \emph{ideals} rather than \emph{elements}. We will show that in a ring of algebraic integers, there is unique factorization, not into prime elements, but into prime ideals. 